\documentclass[12pt]{article}

\usepackage{sbc-template}
\usepackage{graphicx,url}
\usepackage[utf8]{inputenc}
\usepackage[brazil]{babel}
%\usepackage[latin1]{inputenc}  

\sloppy

\title{Ferramenta em Python para Simulação de Dispositivos IoHT, Consenso e Agregação em Federated Learning}
\author{Adriano Zavareze Righi\inst{1}}


\address{Universidade do Vale do Rio dos Sinos
  (UNISINOS)\\
  São Leopoldo -- RS -- Brazil
  \email{contato@adrianorighi.com}
}

\begin{document}

\maketitle

\begin{resumo}
O avanço do \textit{Internet of Health Things} (IoHT) tem impulsionado o uso de dispositivos conectados para monitoramento contínuo da saúde e apoio a decisões clínicas, ao mesmo tempo em que reforça preocupações com privacidade e segurança dos dados sensíveis de pacientes \cite{Kumar2020_SecureIoHT}. Abordagens de \textit{Federated Learning} (FL) surgem como alternativa promissora ao treinamento centralizado, permitindo a construção de modelos colaborativos sem compartilhamento direto de dados, aspecto particularmente relevante em cenários médicos distribuídos \cite{Nevrataki2023_SurveyFLApplications}. Complementarmente, técnicas de consenso e estratégias de agregação de modelos influenciam diretamente a robustez e o desempenho do FL em ambientes com dispositivos heterogêneos e comunicação limitada, como na IoHT \cite{Zhang2024_FLIoTSurvey}. Esse trabalho propõe o desenvolvimento de uma ferramenta em Python para simulação de cenários IoHT com múltiplos dispositivos, suportando diferentes mecanismos de consenso e políticas de agregação em FL, de forma reproduzível e extensível para experimentação acadêmica e avaliação de novas estratégias. A ferramenta visa facilitar o estudo de trade-offs entre acurácia, custo de comunicação, resiliência a falhas e privacidade, utilizando casos de teste inspirados em aplicações de monitoramento de sinais vitais e detecção de anomalias em dados de saúde distribuídos.
\end{resumo}

\section{Introdução}

O conceito de \textit{Internet of Health Things} (IoHT) descreve ecossistemas de dispositivos, sensores e plataformas conectadas voltados à coleta contínua de dados clínicos e ao suporte a serviços inteligentes de saúde \cite{Kumar2020_SecureIoHT}. Nesse contexto, dispositivos vestíveis, sensores ambientais e equipamentos médicos hospitalares geram fluxos de dados distribuídos e sensíveis, cuja centralização em um único repositório pode ser inviável ou indesejável por razões regulatórias, de privacidade e de confiabilidade \cite{Kumar2020_SecureIoHT,Nevrataki2023_SurveyFLApplications}. 

\textit{Federated Learning} (FL) surge como paradigma para treinamento colaborativo de modelos de aprendizado de máquina em que apenas atualizações de parâmetros ou gradientes são compartilhadas com um orquestrador, preservando os dados em seus locais de origem \cite{Zhang2024_FLIoTSurvey}. Em IoHT, o uso de FL é particularmente atrativo para aplicações de monitoramento remoto, diagnóstico assistido e detecção precoce de eventos clínicos, permitindo a colaboração entre instituições e dispositivos sem exposição direta de dados de pacientes \cite{Nevrataki2023_SurveyFLApplications,Kalaiselvi2023_HealthcareMetaverseSurvey}.

Ao mesmo tempo, a heterogeneidade dos dispositivos IoHT, variações de conectividade e riscos de ataques adversariais tornam o problema de consenso e de agregação de modelos um elemento central na eficácia de soluções de FL em saúde \cite{Kumar2020_SecureIoHT,Nevrataki2023_SurveyFLApplications}. Estratégias clássicas como FedAvg podem ser insuficientes em cenários de dados não identicamente distribuídos, presença de clientes maliciosos ou flutuações de participação \cite{Nevrataki2023_SurveyFLApplications,Zhang2024_FLIoTSurvey}. Nesse cenário, ferramentas de simulação em Python têm sido utilizadas para avaliar algoritmos e protocolos de FL em ambientes controlados, permitindo a reprodução de diferentes condições de rede, padrões de participação e políticas de agregação \cite{Intel2020_OpenFLTutorial}.

Esse trabalho apresenta a proposta de desenvolvimento de uma ferramenta de simulação voltada especificamente para cenários de IoHT, com foco em mecanismos de consenso e agregação de FL implementados em Python.

\section{Descrição do Projeto}

\subsection{Visão geral}

A ferramenta proposta consiste em um simulador modular que representa um cenário IoHT composto por múltiplos \textit{nós} correspondentes a dispositivos ou gateways de saúde, um servidor de coordenação de FL e canais de comunicação configuráveis, permitindo a experimentação de diversos algoritmos de consenso e esquemas de agregação de modelos. O objetivo principal é oferecer uma plataforma de estudo que facilite a comparação entre diferentes estratégias de FL sob condições realistas de IoHT, como dados não IID, falhas de comunicação e recursos computacionais limitados \cite{Zhang2024_FLIoTSurvey,Kalaiselvi2023_HealthcareMetaverseSurvey}.

Cada nó IoHT simulado possui um subconjunto de dados representando sinais vitais, medições de sensores ou registros clínicos simplificados, com possibilidade de configurar o grau de heterogeneidade entre dispositivos. O servidor central implementa o ciclo de FL (seleção de clientes, distribuição do modelo global, agregação das atualizações locais) e disponibiliza interfaces para plugar diferentes funções de agregação, desde FedAvg até variantes ponderadas, robustas ou baseadas em valor de contribuição \cite{Zhao2022_FLIoTDataValue}. Além disso, a ferramenta permitirá incluir camadas simples de mecanismos de consenso (por exemplo, votação entre agregadores, checagem de consistência de modelos, filtros de outliers) aplicados sobre os modelos recebidos antes da atualização global.

\subsection{Objetivos}

Entre os objetivos específicos da ferramenta, destacam-se:

\begin{itemize}
    \item Permitir a simulação de cenários IoHT com dezenas a centenas de dispositivos virtuais, com controle sobre distribuição dos dados, padrões de participação e capacidade computacional.
    \item Implementar uma coleção inicial de algoritmos de agregação de FL (ex. FedAvg, agregação ponderada por volume de dados, agregação baseada em valor de contribuição do cliente) e módulos de consenso simples, possibilitando a comparação experimental entre eles \cite{Zhao2022_FLIoTDataValue,Nevrataki2023_SurveyFLApplications}.
    \item Fornecer uma API em Python para que outros pesquisadores possam estender a ferramenta com novos algoritmos de agregação, métodos de consenso, tipos de modelos de aprendizado de máquina e cenários de rede.
    \item Facilitar a definição de casos de teste reprodutíveis, incluindo geração de dados sintéticos inspirados em sinais vitais e métricas de desempenho como acurácia, convergência, custo de comunicação e robustez a falhas ou clientes adversariais.
\end{itemize}

\subsection{Escopo}

Nesse trabalho, o escopo inicial do projeto concentra-se em:

\begin{itemize}
    \item Simulação \textit{offline} (sem integração em tempo real com dispositivos físicos), com dados sintéticos gerados localmente ou conjuntos de dados públicos simplificados.
    \item Implementação de um ciclo básico de FL com modelo de classificação em tarefas relacionadas a saúde (por exemplo, classificação de risco ou detecção de anomalias em sinais vitais simulados).
    \item Suporte a alguns esquemas de agregação e consenso selecionados na literatura recente, deixando a integração com mecanismos avançados (ex., blockchain, privacidade diferencial) como possibilidades para trabalhos futuros \cite{Kumar2020_SecureIoHT,Nevrataki2023_SurveyFLApplications}.
\end{itemize}

\section{Definições e Decisões de Projeto}

\subsection{Arquitetura Proposta}

A arquitetura proposta segue um padrão cliente-servidor típico de FL, estendido com camadas de abstração para cenários IoHT, conforme ilustrado na Figura~\ref{fig:visao-arquitetural}. Em alto nível, a ferramenta será organizada em quatro módulos principais:

\begin{figure}[!ht]
	\caption{Visão arquitetural}
	\label{fig:visao-arquitetural}
	\centering%
	\begin{minipage}{1\textwidth}
		\includegraphics[width=\textwidth]{imagens/visao-arquitetural.png}
	\end{minipage}
\end{figure}

\begin{itemize}
    \item \textbf{Módulo de Orquestração de FL:} responsável pelo gerenciamento das rodadas de treinamento federado, seleção de clientes, distribuição do modelo global, agregação das atualizações e registro de métricas.
    
    \item \textbf{Módulo de Simulação de Dispositivos IoHT:} responsável pela representação de dispositivos e gateways, gerenciamento dos dados locais, treinamento dos modelos locais e simulação de restrições de recursos, como limitação de épocas e latência de comunicação.
    
    \item \textbf{Módulo de Consenso e Agregação:} responsável por encapsular diferentes algoritmos de agregação de modelos e mecanismos leves de consenso aplicados às atualizações recebidas, permitindo a seleção da estratégia em tempo de configuração ou por meio de scripts.
    
    \item \textbf{Módulo de Experimentos e Métricas:} responsável por disponibilizar utilitários para definição de cenários de teste, coleta de métricas, geração de gráficos e exportação de resultados para análise posterior.
\end{itemize}

Do ponto de vista de implementação, a arquitetura seguirá uma abordagem orientada a objetos em Python, com classes representando o servidor de FL, os clientes IoHT e as estratégias de agregação e consenso, conforme ilustrado na Figura~\ref{fig:diagrama-classes}. Essa abordagem busca favorecer a modularidade, a separação de responsabilidades e a extensibilidade da ferramenta.

\begin{figure}[!ht]
	\caption{Diagrama de Classes}
	\label{fig:diagrama-classes}
	\centering%
	\begin{minipage}{1\textwidth}
		\includegraphics[width=\textwidth]{imagens/diagrama-classes.png}
	\end{minipage}
\end{figure}

\subsection{Componentes Principais}

Os principais componentes previstos para a primeira versão da ferramenta são:

\begin{itemize}
    \item \texttt{IoHTClient}: classe responsável por abstrair um dispositivo ou nó IoHT, contendo o subconjunto de dados locais, o modelo local, os parâmetros de treinamento e a lógica de atualização.
    
    \item \texttt{FLServer}: classe responsável por coordenar as rodadas de FL, armazenar o modelo global, selecionar clientes e aplicar a estratégia de agregação.
    
    \item \texttt{Aggregator}: interface ou classe base para algoritmos de agregação. As implementações concretas incluirão FedAvg, variantes ponderadas e versões que considerem o valor de contribuição dos dados, conforme proposto em \cite{Zhao2022_FLIoTDataValue}.
    
    \item \texttt{ConsensusModule}: componente opcional responsável por aplicar mecanismos simplificados de consenso ou validação sobre os modelos enviados pelos clientes, auxiliando na detecção de comportamentos anômalos ou inconsistentes \cite{Kumar2020_SecureIoHT}.
    
    \item \texttt{ExperimentRunner}: camada responsável pela orquestração de experimentos completos, desde a geração ou carregamento dos dados até a execução de múltiplas rodadas de FL com diferentes configurações.
\end{itemize}

O fluxo geral de execução previsto, representado na Figura~\ref{fig:diagrama-fluxo} consiste resumidamente em: geração ou carregamento dos dados, distribuição entre os clientes IoHT, treinamento local dos modelos, aplicação das etapas de validação ou consenso, agregação global dos parâmetros e coleta das métricas experimentais.

\begin{figure}[!ht]
	\caption{Diagrama de Fluxo}
	\label{fig:diagrama-fluxo}
	\centering%
	\begin{minipage}{1\textwidth}
		\includegraphics[width=\textwidth]{imagens/diagrama-fluxo.png}
	\end{minipage}
\end{figure}

\subsection{Tecnologias e Ferramentas}

A implementação será realizada em Python, aproveitando o ecossistema de bibliotecas amplamente adotadas em pesquisas relacionadas a FL e IoHT:

\begin{itemize}
    \item \textbf{Linguagem:} Python 3.x, devido à ampla disponibilidade de ferramentas de aprendizado de máquina e suporte à simulação de processos distribuídos \cite{Intel2020_OpenFLTutorial}.
    
    \item \textbf{Bibliotecas de ML:} frameworks como PyTorch ou TensorFlow serão avaliados. Entretanto, para o escopo da disciplina, a tendência é utilizar PyTorch devido à sua simplicidade de integração em cenários de simulação de FL \cite{Intel2020_OpenFLTutorial}.
    
    \item \textbf{Bibliotecas de apoio:} NumPy e Pandas para manipulação de dados, Matplotlib ou Seaborn para visualização de métricas, além de bibliotecas voltadas à simulação de eventos ou \textit{multi-threading} leve, conforme necessidade.
    
    \item \textbf{Ferramentas de FL:} bibliotecas de código aberto, como OpenFL, serão utilizadas como referência arquitetural e de boas práticas para federated learning aplicado à saúde, embora o simulador proposto seja desenvolvido de forma independente e simplificada sendo talvez desnecessário seu uso, com foco em extensibilidade acadêmica e experimentação \cite{Intel2020_OpenFLTutorial}.
\end{itemize}

Além disso, o código-fonte será versionado em um repositório Git e disponibilizado em plataforma pública para análise, reprodução e reuso pela comunidade acadêmica.

\subsection{Decisões de Projeto}

Algumas decisões preliminares de projeto incluem:

\begin{itemize}
    \item \textbf{Modelo de dados:} utilização de dados sintéticos ou conjuntos públicos simplificados, com foco em tarefas de classificação binária ou multiclasse relacionadas a sinais vitais ou riscos à saúde, permitindo rápida experimentação e reduzindo barreiras de acesso a dados reais.
    
    \item \textbf{Topologia de rede:} adoção inicial de uma topologia em estrela, na qual um servidor central coordena os clientes, refletindo o padrão mais comum em FL aplicado a IoHT \cite{Zhang2024_FLIoTSurvey}.
    
    \item \textbf{Consenso:} implementação de mecanismos simplificados de consenso, como verificação de desvio em relação à mediana dos parâmetros e rejeição de modelos fora de um intervalo esperado, funcionando como etapa de validação antes da agregação global, com possibilidade de extensão futura.
    
    \item \textbf{Extensibilidade:} definição de interfaces claras para agregadores e módulos de consenso, permitindo a introdução de novos algoritmos por meio da criação de subclasses, sem necessidade de alteração do núcleo da ferramenta.
\end{itemize}

O foco principal da proposta está na avaliação experimental de cenários de FL aplicados a IoHT, priorizando simplicidade de implementação, modularidade e facilidade de extensão, em vez de requisitos de escalabilidade de infraestrutura em nível de produção.

\section{Metodologia de Avaliação e Resultados}
\label{sec:avaliacao}

\subsection{Cenários experimentais}

Foram definidos cinco cenários para demonstrar a ferramenta e avaliar o impacto de diferentes estratégias de agregação e consenso:

\begin{table}[!ht]
\centering
\caption{Parâmetros dos cenários experimentais.}\label{tab:cenarios}
\begin{tabular}{lccccc}
\hline
\textbf{Cenário} & \textbf{Agregador} & \textbf{Consenso} & \textbf{Clientes} & \textbf{Rodadas} & \textbf{non-IID} \\
\hline
Baseline    & FedAvg    & Threshold & 50 & 15 & Não \\
Non-IID     & Weighted  & Voting    & 50 & 15 & Sim \\
Robusto     & Robust    & Threshold & 50 & 15 & Sim \\
MultiKrum   & MultiKrum & Threshold & 20 & 10 & Não \\
HotStuff    & FedAvg    & HotStuff  & 20 & 10 & Não \\
\hline
\end{tabular}
\end{table}

O cenário \textbf{baseline} representa a configuração padrão de referência: FedAvg com \textit{ThresholdConsensus}, dados aproximadamente IID (sem drift), taxa de falha de 5\% e ruído moderado de 3\%. O cenário \textbf{non-IID} utiliza \textit{WeightedAggregator} com \textit{VotingConsensus} (concordância mínima de 66\%), distribuição não IID entre clientes com drift configurável, taxa de falha de 8\% e ruído de 4\%. O cenário \textbf{robusto} combina \textit{RobustAggregator} com \textit{ThresholdConsensus} (desvio máximo de 0,5), dados não IID, taxa de falha de 10\% e ruído de 8\%, representando condições mais adversas.

Os cenários \textbf{MultiKrum} e \textbf{HotStuff} foram adicionados posteriormente para explorar estratégias avançadas: \textit{MultiKrumAggregator} com tolerância a $f=2$ bizantinos, e \textit{HotStuffConsensus} com comitê BFT ($f=2$, comitê de 7 validadores). Ambos utilizam 20 clientes, 100 amostras por cliente e 10 rodadas, com taxa de falha de 10\%.

Todos os cenários utilizam fração de participação de 75\% por rodada e semente aleatória fixa (\texttt{seed = 42}) para reprodutibilidade. A ferramenta gera arquivos CSV e JSON com métricas por rodada na pasta \texttt{results/}, permitindo análise posterior.

\subsection{Métricas coletadas}

As principais métricas avaliadas foram:
\begin{itemize}
  \item \textbf{Participação}: número de clientes selecionados por rodada.
  \item \textbf{Atualizações aceitas}: número de atualizações que passaram pelo mecanismo de consenso.
  \item \textbf{Custo de comunicação}: proxy proporcional ao número de clientes participantes em cada rodada.
  \item \textbf{Proxy de acurácia}: medida heurística baseada na magnitude dos parâmetros do modelo global, variando entre 0,5 e 0,99.
\end{itemize}

Embora o proxy de acurácia não corresponda a uma métrica real de desempenho em dados de teste, ele permite observar tendências relativas entre cenários e a capacidade da ferramenta de convergir para parâmetros mais estáveis ao longo das rodadas.

\subsection{Resultados e discussão}

A Tabela~\ref{tab:resultados} consolida as métricas obtidas para cada cenário, incluindo valores da primeira e última rodada, além da rodada em que o proxy de acurácia atingiu 0,99 (saturação).

\begin{table}[!ht]
\centering
\caption{Métricas consolidadas por cenário experimental.}\label{tab:resultados}
\begin{tabular}{lcccccc}
\hline
\textbf{Cenário} & \textbf{Partic.} & \textbf{Aceitação} & \textbf{Proxy R1} & \textbf{Proxy final} & \textbf{Saturação} & \textbf{Custo} \\
\hline
Baseline    & 37/50 & 33--37 & 0,537 & 0,990 & Rod. 14 & 111 \\
Non-IID     & 37/50 & 32--37 & 0,546 & 0,990 & Rod. 11 & 111 \\
Robusto     & 37/50 & 30--36 & 0,548 & 0,990 & Rod. 12 & 111 \\
MultiKrum   & 15/20 & 13--15 & 0,535 & 0,868 & --- & 45 \\
HotStuff    & 15/20 & 7      & 0,538 & 0,892 & --- & 45 \\
\hline
\end{tabular}
\end{table}

\textbf{Cenário baseline (FedAvg + Threshold).} A configuração de referência apresentou crescimento estável do proxy de acurácia, partindo de 0,537 na rodada 1 e atingindo 0,99 na rodada 14. A taxa de aceitação de updates foi elevada (33--37 de 37 participantes por rodada), refletindo a baixa taxa de falha (5\%) e o limiar de desvio tolerante (\texttt{max\_deviation = 0,6}). Os parâmetros do modelo global evoluíram de forma monotônica ($w_1$: 0,08 $\to$ 1,14; $w_2$: 0,09 $\to$ 1,39; \textit{bias}: 0,05 $\to$ 0,73), indicando convergência consistente.

\textbf{Cenário non-IID (Weighted + Voting).} Apesar da distribuição não IID dos dados, este cenário atingiu saturação do proxy de acurácia na rodada 11 -- três rodadas antes do baseline. Esse resultado sugere que a ponderação adicional do \textit{WeightedAggregator} (baseada no valor médio das saídas de cada cliente) compensou a heterogeneidade dos dados, acelerando a convergência. A taxa de aceitação do \textit{VotingConsensus} (concordância mínima de 66\%) foi similar à do baseline (32--37 de 37), indicando que o critério de moda por coeficiente não foi excessivamente restritivo. Os parâmetros finais apresentaram magnitudes ligeiramente superiores ($w_1$: 1,25; $w_2$: 1,55; \textit{bias}: 1,24), coerentes com a ponderação mais agressiva.

\textbf{Cenário robusto (Robust + Threshold).} Este cenário combinou a maior taxa de falha (10\%) com o limiar de desvio mais restritivo (\texttt{max\_deviation = 0,5}) e ruído elevado (8\%). A taxa de aceitação variou entre 30 e 36 de 37 participantes, a maior variância entre os cenários, refletindo as condições adversas. O proxy de acurácia atingiu 0,99 na rodada 12, demonstrando que o \textit{RobustAggregator} (baseado na mediana dos parâmetros) mitigou efetivamente o impacto de atualizações extremas e falhas, mesmo com ruído mais intenso.

\textbf{Cenário MultiKrum (MultiKrum + Threshold).} Com 20 clientes e 10 rodadas, o MultiKrum apresentou crescimento mais conservador do proxy de acurácia (0,535 $\to$ 0,868), sem atingir saturação. Esse comportamento é esperado, dado que o MultiKrum descarta as $f$ atualizações mais divergentes antes da agregação, reduzindo o número efetivo de contribuições por rodada. A aceitação manteve-se entre 13 e 15 de 15 participantes, indicando que o \textit{ThresholdConsensus} atuou de forma complementar na filtragem de outliers.

\textbf{Cenário HotStuff (FedAvg + HotStuff).} O \textit{HotStuffConsensus} opera por comitê BFT: com $f = 2$, o comitê tem tamanho $3f + 1 = 7$. Assim, apenas 7 atualizações foram aceitas por rodada, de forma consistente, independentemente do número de participantes (15). Apesar da redução drástica de participantes efetivos, o proxy de acurácia evoluiu de 0,538 para 0,892 ao longo de 10 rodadas, evidenciando que o consenso baseado em comitê consegue manter a qualidade da agregação mesmo com participação restrita. Esse comportamento ilustra o \textit{trade-off} entre segurança BFT e eficiência de aprendizado.

Em síntese, os resultados confirmam que diferentes combinações de agregadores e mecanismos de consenso produzem trajetórias de convergência distintas, mesmo sobre um proxy simplificado de acurácia. O \textit{ThresholdConsensus} mostrou-se eficaz e de baixo custo computacional, filtrando atualizações com desvio excessivo em todos os cenários. O \textit{VotingConsensus} introduziu um critério alternativo de concordância que, combinado com ponderação por valor de dados, acelerou a convergência em dados não IID. O \textit{HotStuffConsensus} demonstrou viabilidade de consenso BFT em FL, com aceitação rigorosa porém estável. Esses resultados estão alinhados com discussões na literatura sobre a importância de estratégias de agregação e consenso em FL aplicado a IoT e saúde .

\section{Conclusão}
\label{sec:conclusao}

Este artigo apresentou o desenvolvimento completo e a demonstração de uma ferramenta em Python para simulação de dispositivos IoHT, consenso e agregação em Federated Learning, em consonância com os requisitos de implementação e avaliação da disciplina de Computação Aplicada do PPGCA \cite{Kumar2020_SecureIoHT,Nevrataki2023_SurveyFLApplications}. A ferramenta foi projetada com arquitetura modular, suportando múltiplas estratégias de agregação e consenso, geração de dados sintéticos, configurações flexíveis de cenários e exportação de métricas para análise.

Como contribuições, destacam-se: (i) a disponibilização de um código-fonte simples e extensível para FL em IoHT, adequado a contextos de ensino; (ii) a integração de diferentes mecanismos de consenso e agregação, incluindo estratégias avançadas como MultiKrum e HotStuff BFT, permitindo estudar sua interação em cenários heterogêneos; e (iii) a demonstração de cinco cenários experimentais representativos, com resultados reprodutíveis e análise quantitativa de métricas por rodada.

Trabalhos futuros incluem a integração de modelos de aprendizado de máquina mais realistas (por exemplo, redes neurais com frameworks como PyTorch), a inclusão de mecanismos de privacidade diferencial e segurança, e a ligação da ferramenta a dispositivos físicos ou emuladores de IoHT. Também se planeja disponibilizar o repositório em plataformas públicas (por exemplo, GitHub) e ampliar a documentação, seguindo recomendações de boas práticas para ferramentas de FL em saúde.

\bibliographystyle{sbc}
\bibliography{references}

\end{document}